\section{Solver Architecture}
\label{sec:solver}

\subsection{Governing equations}

We solve the incompressible Navier--Stokes equations,
\begin{align}
    \nabla \cdot \bm{u} &= 0, \label{eq:continuity} \\
    \frac{\partial \bm{u}}{\partial t} + (\bm{u} \cdot \nabla)\bm{u} &= -\nabla p + \nu \nabla^2 \bm{u} + \bm{f}, \label{eq:momentum}
\end{align}
where $\bm{u}$ is the velocity field, $p$ the kinematic pressure, $\nu$ the kinematic viscosity, and $\bm{f}$ a body force (e.g., a mean pressure gradient driving channel flow).
When a turbulence closure is active, the molecular viscosity $\nu$ is replaced by an effective viscosity $\nueff = \nu + \nut$, where $\nut$ is the eddy viscosity predicted by the closure model.
For anisotropic closures (EARSM, TBNN, TBRF), the Reynolds stress divergence is added directly as a source term rather than through an eddy viscosity.

\subsection{Fractional-step projection method}

Time integration follows the fractional-step (pressure-projection) method:
\begin{enumerate}
    \item \textbf{Turbulence update}: Advance the turbulence model to compute updated $\nut$ (or anisotropic stress tensor $\tau_{ij}$).
    \item \textbf{Predictor}: Compute an intermediate velocity $\bm{u}^*$ by advancing the momentum equation without the pressure gradient:
    \begin{equation}
        \bm{u}^* = \bm{u}^n + \Delta t \left( -\bm{C}(\bm{u}^n) + \bm{D}(\bm{u}^n) + \bm{f} \right),
    \end{equation}
    where $\bm{C}$ and $\bm{D}$ denote the convective and diffusive operators, respectively.
    Convection uses the skew-symmetric form for improved energy conservation.
    \item \textbf{Pressure Poisson}: Solve
    \begin{equation}
        \nabla^2 p' = \frac{1}{\Delta t} \nabla \cdot \bm{u}^*
    \end{equation}
    for the pressure correction $p'$.
    \item \textbf{Velocity correction}: Project to a divergence-free field:
    \begin{equation}
        \bm{u}^{n+1} = \bm{u}^* - \Delta t \, \nabla p'.
    \end{equation}
\end{enumerate}
This sequence is embedded within a strong stability preserving Runge--Kutta (SSP-RK3) time integrator, with steps 2--4 repeated at each stage with appropriate SSP weights.
Adaptive time stepping enforces directional CFL constraints: a relaxed $\mathrm{CFL}_{xz}$ for the streamwise and spanwise directions (where grid spacing is uniform) and a strict $\mathrm{CFL}_{\max}$ for the wall-normal direction (where the stretched grid produces small $\Delta y$ near walls).

\subsection{Staggered MAC grid}

Spatial discretization uses a staggered marker-and-cell (MAC) grid with second-order central differences.
Velocity components are stored at cell faces ($u$ at $x$-faces, $v$ at $y$-faces, $w$ at $z$-faces) and pressure at cell centers.
This arrangement guarantees pressure--velocity coupling without checkerboard oscillations and ensures that the discrete divergence and gradient operators compose to the exact discrete Laplacian:
\begin{equation}
    D \cdot G = L,
\end{equation}
a property critical for the projection method to enforce the incompressibility constraint~\eqref{eq:continuity} to machine precision.
For stretched grids in the wall-normal direction, this compatibility is maintained through precomputed metric arrays ($\Delta y_v$, $\Delta y_c$, $y_{\mathrm{Lap},*}$) that encode the non-uniform spacing in the discrete operators.

\subsection{GPU acceleration}

The solver achieves GPU acceleration through OpenMP target offload directives, maintaining a single source code path for both CPU and GPU execution.
All field data (velocity, pressure, turbulence quantities) is persistently mapped to GPU memory at initialization and remains on device throughout the simulation.
GPU--CPU data transfer occurs only during I/O operations (writing output files) and at initialization/finalization.

Key design decisions for GPU performance include:
\begin{itemize}
    \item \textbf{Single code path}: Arithmetic kernels contain no \texttt{\#ifdef} guards---the OpenMP target pragmas are silently ignored in CPU builds, ensuring identical numerical behavior on both platforms.
    \item \textbf{Free-function kernels}: GPU compute kernels are implemented as free functions rather than class methods to avoid implicit transfer of the \texttt{this} pointer (an nvc++ compiler requirement).
    \item \textbf{Persistent device memory}: The \texttt{OmpDeviceBuffer} wrapper manages GPU allocations with RAII semantics. Data is mapped once and remains on device.
    \item \textbf{Reduction-based diagnostics}: Scalar diagnostics (kinetic energy, maximum velocity, divergence $L_\infty$ norm) are computed via GPU reductions without transferring full fields to the host.
\end{itemize}

The solver is compiled with the NVHPC compiler (nvc++) for GPU builds and GCC for CPU-only builds.

\subsection{Poisson solver}

The pressure Poisson equation is solved using a fast Fourier transform (FFT) approach via the cuFFT library.
For directions with periodic boundary conditions, standard FFTs are applied; for directions with Dirichlet or Neumann conditions, discrete cosine/sine transforms are used.
The FFT approach reduces the 3D Poisson problem to a set of independent tridiagonal systems in the wall-normal direction, which are solved directly.
For stretched grids, the tridiagonal coefficients incorporate the mesh metric arrays to maintain $D \cdot G = L$ consistency, and the right-hand side uses volume-weighted mean subtraction to satisfy solvability.

The FFT Poisson solver achieves 22--32$\times$ speedup over the geometric multigrid alternative on GPU for stretched grids, making it the preferred backend for all production runs.

\subsection{Immersed boundary method}

Complex geometries (cylinder, airfoil, periodic hills) are handled via an immersed boundary method (IBM) using direct forcing.
At initialization, signed distance functions define the solid boundary, and weight arrays are precomputed: $w = 0$ inside the solid, $w = 1$ in the fluid, with a smooth transition in the forcing region.
During each time step, the IBM applies element-wise multiplication of the velocity field by these weight arrays---a trivially parallel operation that adds less than 0.3\% overhead on GPU.
The Poisson equation right-hand side is masked at solid cells via a separate GPU kernel, avoiding any CPU--GPU synchronization.

\subsubsection{Ghost-cell immersed boundary method}
\label{sec:ghost_cell_ibm}

The primary IBM approach uses the ghost-cell method of \citet{Fadlun2000}.
For each ghost cell (the first solid cell adjacent to the fluid--solid interface), the velocity is reconstructed by interpolating from the fluid side along the surface-normal direction to enforce the no-slip condition at the boundary.
This yields sharp boundary representation with second-order accuracy in the boundary location, avoiding the $O(\eta)$ accuracy limitation of volume penalization methods~\citep{Angot1999}.

The solver also retains a semi-implicit volume penalization option,
\begin{equation}
    \bm{u} \leftarrow \bm{u} \left[ 1 - (1 - w) \min\!\left(\frac{\Delta t}{\eta}, 1\right) \right],
    \label{eq:volume_penalization}
\end{equation}
where $\eta$ is a penalization timescale and $w$ is the solid mask ($w=0$ inside the body, $w=1$ in the fluid).
This penalization is used during the SST warm-up phase for the periodic hills geometry, where it provides robust stabilization of the initial transient; the solver then switches to the ghost-cell method for the production measurement phase.
For cylinder and sphere cases, ghost-cell forcing is used throughout.

In all cases, the re-forcing step after the pressure correction is removed: only the pre-correction forcing (step~5 in Section~\ref{sec:solver}) and the Poisson RHS masking (step~7) are retained.
This eliminates a feedback instability that arises when the corrected velocity is re-forced, generating divergence that is never projected out.

\subsubsection{Grid convergence of the ghost-cell IBM}
\label{sec:ibm_grid_convergence}

The ghost-cell method of \citet{Fadlun2000} is first-order accurate in the grid spacing $h$, since the velocity at each ghost cell is reconstructed via linear interpolation along the surface normal.
We verify this convergence behavior on two canonical cases: the circular cylinder at $\Rey = 100$ and the sphere at $\Rey = 200$.

Table~\ref{tab:ibm_convergence_cyl} shows the cylinder drag coefficient $C_D$ as a function of grid resolution, expressed as cells per diameter $D/h$.
At 8 cells/$D$ the error exceeds 50\%, but by 24 cells/$D$ the ghost-cell method recovers $C_D$ to within 9\% of the reference value $C_D = 1.35$.
Adding SST warm-up (Section~\ref{sec:sst_warmup}) further improves the result to $C_D = 1.40$ (3.6\% error), indicating that the turbulence initialization contributes meaningfully once the IBM resolution is adequate.

\begin{table}[htbp]
\centering
\caption{Grid convergence of ghost-cell IBM for circular cylinder at $\Rey = 100$. Reference $C_D = 1.35$.}
\label{tab:ibm_convergence_cyl}
\begin{tabular}{lrrr}
\toprule
Grid & Cells/$D$ & $C_D$ & Error \\
\midrule
$128 \times 96$  & 8  & 0.63 & 53\% \\
$384 \times 288$ & 24 & 1.47 & 9\%  \\
$384 \times 288$ (SST warm-up) & 24 & 1.40 & 3.6\% \\
\bottomrule
\end{tabular}
\end{table}

Table~\ref{tab:ibm_convergence_sph} presents the corresponding data for the sphere.
Convergence is substantially slower: at 24 cells/$D$, the error remains 51\%.
This is a known property of first-order immersed boundary methods on staggered grids---the three-dimensional interpolation geometry at curved surfaces is more complex than in two dimensions, and the volume fraction of cells cut by the interface scales as $O(h^{-2})$ in 3D versus $O(h^{-1})$ in 2D.
Notably, the baseline simulation (no SST warm-up) produces the same $C_D$ as the SST-initialized run at $384 \times 256^2$, confirming that the IBM grid resolution is the convergence bottleneck rather than the turbulence model.

\begin{table}[htbp]
\centering
\caption{Grid convergence of ghost-cell IBM for sphere at $\Rey = 200$. Reference $C_D = 0.77$.}
\label{tab:ibm_convergence_sph}
\begin{tabular}{lrrr}
\toprule
Grid & Cells/$D$ & $C_D$ & Error \\
\midrule
$128 \times 96^2$  & 10 & 0.17 & 78\% \\
$192 \times 128^2$ & 13 & 0.25 & 68\% \\
$256 \times 192^2$ & 17 & 0.21 & 73\% \\
$384 \times 256^2$ & 24 & 0.38 & 51\% \\
\bottomrule
\end{tabular}
\end{table}

These results are consistent with the $O(h)$ convergence rate reported in the IBM literature \citep{Fadlun2000}.
For the cylinder, 24 cells/$D$ provides sufficient resolution for quantitative $C_D$ comparison (${<}10\%$ error).
For the sphere, while absolute $C_D$ values have not converged at affordable resolutions, relative comparisons between turbulence models at a fixed grid resolution remain meaningful---any model-to-model differences are attributable to the closure rather than to IBM discretization error.

\subsubsection{Grid resolution requirements for IBM stability}
\label{sec:ibm_grid_requirements}

The combination of the explicit RK3 time integrator with the ghost-cell IBM imposes minimum grid resolution requirements that are substantially higher than those for body-fitted solvers.
The IBM introduces a pressure transient at the fluid--solid interface during the initial time steps: the penalization term damps the body-interior velocity toward zero, creating a localized divergence source that the Poisson solver must correct.
With insufficient resolution, this transient produces oscillations that grow exponentially and diverge within a few hundred steps.

Table~\ref{tab:ibm_min_grid} reports the minimum stable grid for each a~posteriori case, determined empirically by testing progressively finer grids until long-time stability is achieved (${>}2000$ steps without divergence).
The periodic hills geometry requires approximately $4\times$ the resolution that would be sufficient for a body-fitted solver, while the sphere requires $2\times$ refinement in each direction.
This overhead is a direct consequence of the explicit time integrator: implicit solvers (e.g., OpenFOAM's SIMPLE algorithm, commonly used for RANS periodic hills at Re${=}5{,}600$) can accommodate the IBM pressure transient without excessive resolution.

\begin{table}[htbp]
\centering
\caption{Minimum stable grid resolution for each IBM case with the explicit RK3 solver. The duct uses wall-fitted boundaries (no IBM). The ``factor'' column indicates the grid refinement relative to an initial estimate based on cell-count matching with body-fitted RANS literature.}
\label{tab:ibm_min_grid}
\begin{tabular}{lrrrr}
\toprule
Case & Grid & Cells & Cells/$D$ & Factor \\
\midrule
Cylinder Re${=}100$ & $384 \times 288$ & 0.1M & 24 & $1\times$ \\
Hills Re${=}5{,}600$ & $1{,}536 \times 768$ & 1.2M & --- & $4\times$ \\
Duct Re$_b{=}3{,}500$ & $96^3$ & 0.9M & --- & $1\times$ \\
Sphere Re${=}200$ & $384 \times 256^2$ & 25.2M & 48 & $2\times$/dir \\
\bottomrule
\end{tabular}
\end{table}

The volume penalization (Eq.~\ref{eq:volume_penalization}) with $\eta = 0.1$ is applied during the SST warm-up phase for all IBM cases to stabilize the initial transient.
After warm-up, the solver switches to ghost-cell forcing for higher accuracy.

\subsection{SST warm-up initialization}
\label{sec:sst_warmup}

Non-transport turbulence models---including mixing-length, GEP, and all neural network closures---compute the eddy viscosity $\nut$ from the local velocity gradient $\partial u_i / \partial x_j$.
When the simulation is initialized from a uniform velocity field, all velocity gradients are zero and these models produce $\nut = 0$, causing the early-time flow development to proceed as if laminar.
This can delay or entirely prevent the establishment of a turbulent mean flow.

To address this, the solver implements a two-phase initialization:
\begin{enumerate}
    \item \textbf{SST warm-up}: The Menter SST $k$--$\omega$ transport model~\citep{Menter1994} is run for a specified physical time (\texttt{warmup\_time}), developing the velocity, $k$, and $\omega$ fields from the uniform initial condition. Because the SST model integrates transport equations for $k$ and $\omega$, it produces nonzero $\nut$ even from uniform flow via the production terms.
    \item \textbf{Model switch}: The target turbulence closure replaces SST and operates on the developed flow field. The $k$ and $\omega$ fields are discarded (unless the target model is itself a transport model).
\end{enumerate}
Timer statistics are reset at the model switch so that reported wall-clock times reflect only the cost of the target closure, not the SST warm-up phase.

\subsection{Three-phase timing methodology}
\label{sec:timing_methodology}

Rigorous cost measurement requires separating model startup transients from steady-state operation.
Each model evaluation uses a three-phase protocol:
\begin{enumerate}
    \item \textbf{SST warm-up} (Section~\ref{sec:sst_warmup}): Develop the flow field using SST $k$--$\omega$ for a specified physical time. This phase is not timed.
    \item \textbf{Model stabilization}: The target closure runs for a specified number of warm-up steps (\texttt{warmup\_steps}) to allow its $\nut$ field to equilibrate with the developed flow. This phase is not timed.
    \item \textbf{Measurement}: Timing is recorded over the remaining steps. Only this phase contributes to the reported wall-clock cost.
\end{enumerate}
All three phases use adaptive time stepping with identical CFL constraints, so the time step adjusts naturally to the effective viscosity of each model.
For the fixed-$\Delta t$ cost comparisons (Section~\ref{sec:cost}), a common time step is imposed across all models to ensure that the Poisson, convection, and diffusion costs are identical and only the turbulence closure cost varies.
The relevant configuration parameters are \texttt{warmup\_model}, \texttt{warmup\_time}, and \texttt{warmup\_steps}.

\subsection{Performance validation}
\label{sec:perf_validation}

To establish that our solver is well-optimized---and therefore that the turbulence closure cost measurements reported in Section~\ref{sec:cost} are meaningful---we benchmark against CaNS \citep{Costa2018} and CaLES \citep{Xiao2024}, widely used GPU-accelerated incompressible Navier--Stokes solvers.
CaNS employs the same numerical approach as our solver: second-order staggered finite differences, SSP-RK3 time integration, and an FFT-based direct Poisson solver.
It is written in Fortran with CUDA Fortran (CUF) kernels and uses the cuDecomp library for multi-GPU pencil decomposition.
CaLES extends CaNS with LES subgrid-scale models and reports single-GPU performance on A100.

We match the CaNS benchmark configuration exactly: channel flow at $\Rey_\tau \approx 590$ on a $512 \times 256 \times 144$ grid (18.9M cells), with second-order central differences, explicit RK3 time integration, explicit diffusion, and FFT-based Poisson solver.

\begin{table}[htbp]
\centering
\caption{Solver throughput comparison with CaNS \citep{Costa2018} and CaLES \citep{Xiao2024} on identical discretization (2nd-order central FD, RK3, explicit diffusion, FFT Poisson). CaNS grid: $512 \times 256 \times 144$ (18.9M cells), V100-SXM2-32GB on DGX-2. CaLES grid: $512 \times 384 \times 1440$ (283M cells), single A100-SXM4-64GB.}
\label{tab:solver_throughput}
\begin{tabular}{llrrr}
\toprule
Solver & GPU configuration & ms/step & Mcells/s & Mcells/s/GPU \\
\midrule
This work   & $1 \times$ V100-PCIe-16GB  & 74.5  & 253  & 253 \\
This work   & $1 \times$ H100-SXM5-80GB  & ---   & ---  & --- \\
\midrule
CaNS        & $4 \times$ V100-SXM2-32GB  & 481    & 39   & 9.8 \\
CaNS        & $16 \times$ V100-SXM2-32GB & 140    & 135  & 8.4 \\
\bottomrule
\end{tabular}
\end{table}

Table~\ref{tab:solver_throughput} shows that our solver achieves 253~Mcells/s on a single V100 at 16.8M cells, which is $6.5\times$ the total throughput of CaNS running on 4~V100s ($39$~Mcells/s) at a comparable grid size, and $1.9\times$ CaNS on 16~V100s ($135$~Mcells/s).
The per-GPU efficiency is $26\times$ higher (253 vs.\ 9.8~Mcells/s/GPU).

The large per-GPU advantage over CaNS's multi-GPU configuration arises from the distributed FFT communication overhead inherent in CaNS's pencil decomposition.
CaNS uses cuDecomp for 2D pencil transposes that require global all-to-all communication.
With RK3, three Poisson solves are performed per time step, each requiring forward and backward FFTs with multiple transposes.
\citet{Costa2018} report that the Poisson solver consumes 60--70\% of wall-clock time in the multi-GPU configuration, with much of that attributable to transposes.
Our single-GPU implementation avoids this overhead entirely, performing the full FFT on-chip via cuFFT.

\subsubsection{Profiling-guided optimization}

To close the remaining gap with CaLES \citep{Xiao2024}---a single-GPU LES code built on CaNS---we used NVIDIA Nsight Systems (nsys) to profile our solver at the kernel level.
The initial profile on V100 at $256^3$ revealed that the Poisson solver consumed 73\% of kernel time, compared to CaLES's 38\%.

The root cause was identified by examining the cuFFT plan configuration.
Our batched 2D FFT used a strided memory layout (\texttt{istride}=$N_y$, \texttt{idist}=1), meaning consecutive elements in the transform dimensions ($x$, $z$) were separated by $N_y = 256$ doubles (2~KB) in memory.
This non-unit stride caused non-coalesced memory access within cuFFT's internal radix kernels: each warp of 32~threads accessed 32 different cache lines instead of 1--2 contiguous lines.
A standalone cuFFT benchmark confirmed an $11\times$ penalty: 8.25~ms/call with stride=$N_y$ versus 0.74~ms/call with stride=1 for identical transforms.

The fix required reorganizing the packed data layout from $\text{packed}[(i N_z + k) N_y + j]$ (y-interleaved, stride=$N_y$ for FFT) to $\text{packed}[j (N_x N_z) + i N_z + k]$ (y-batched, stride=1 for FFT).
Since the cuSPARSE batched tridiagonal solver requires mode-major layout ($\text{hat}[m N_y + j]$), we added a shared-memory tiled matrix transpose (32$\times$32 tiles with bank-conflict-free padding) before and after the tridiagonal solve.
The transpose cost ($\sim$0.3~ms/call) is negligible compared to the cuFFT savings.

A second optimization separated the post-FFT data unpacking and boundary condition application from a single kernel with scattered writes into two kernels: a pure transpose (output-coalesced, $4.7\times$ faster) and a separate ghost cell fill (no branch divergence).

Table~\ref{tab:optimization} summarizes the cumulative impact.

\begin{table}[htbp]
\centering
\caption{Profiling-guided optimizations and their impact on total step time (V100, $256^3$, channel flow, RK3).}
\label{tab:optimization}
\begin{tabular}{lrrr}
\toprule
Version & ms/step & ns/point & Cumulative speedup \\
\midrule
Initial implementation                 & 141.0 & 8.40 & $1.00\times$ \\
+ Output-coalesced unpack kernel       & 123.9 & 7.39 & $1.14\times$ \\
+ Contiguous cuFFT layout (stride=1)   & 74.5  & 4.44 & $1.89\times$ \\
\bottomrule
\end{tabular}
\end{table}

After optimization, the kernel-level breakdown (nsys) shows a well-balanced solver with no single dominant bottleneck (Table~\ref{tab:kernel_breakdown}).

\begin{table}[htbp]
\centering
\caption{GPU kernel time breakdown after optimization (V100, $256^3$, 20 steps). Percentages are of total kernel time.}
\label{tab:kernel_breakdown}
\begin{tabular}{lrr}
\toprule
Category & ms/step & \% of kernels \\
\midrule
Poisson solver (FFT + tridiag + pack/unpack) & 35.0 & 30\% \\
\quad cuFFT transforms                       & 16.6 & 14\% \\
\quad cuSPARSE tridiagonal                   & 10.5 & 9\% \\
\quad Pack/unpack + transpose                & 7.9  & 7\% \\
RK3 stepping (blend + substep)               & 33.5 & 28\% \\
Convection + diffusion                       & 17.3 & 15\% \\
Velocity correction + projection             & 18.3 & 16\% \\
Divergence                                   & 6.8  & 6\% \\
Boundary conditions                          & 2.3  & 2\% \\
Other                                        & 3.5  & 3\% \\
\bottomrule
\end{tabular}
\end{table}

We also evaluated replacing the cuSPARSE PCR-based batched tridiagonal solver with a custom Thomas algorithm (as used by CaNS), but found cuSPARSE to be faster at $N_y = 256$: the PCR algorithm's $O(\log N)$ parallel steps outperform Thomas's $O(N)$ serial sweep when sufficient inter-system parallelism exists, and the Thomas solver's per-thread local array ($N_y$ doubles) causes register spilling to local memory on GPU.

\subsubsection{Comparison with CaLES}

CaLES \citep{Xiao2024} reports a per-subroutine timing breakdown for its base solver (DNS mode, no turbulence model) on a single A100-SXM4-64GB GPU with a $512 \times 384 \times 1440$ grid (283M cells):
Poisson solver 38\%, RHS (convection + diffusion) 30\%, velocity correction 13\%, and SGS model 19\% (Smagorinsky, not applicable to our RANS comparison).
Their DNS throughput is 1.50~ns/step/point ($\sim$667~Mcells/s).

Our optimized solver achieves 4.44~ns/step/point on V100 at $256^3$.
Accounting for the V100/A100 memory bandwidth ratio (0.90 vs.\ 2.0~TB/s), the bandwidth-normalized per-point cost is within a factor of 1.3 of CaLES.
The residual gap is attributable to (i) smaller grid size (16.8M vs.\ 283M cells, reducing GPU occupancy), (ii) programming model differences (OpenMP target offload vs.\ CUDA Fortran/OpenACC), and (iii) CaNS's in-place cuFFT execution which avoids one buffer allocation.

These results confirm that our solver is competitive with established, heavily optimized GPU codes, and that the turbulence closure cost measurements reported in Section~\ref{sec:cost} are measured against a well-optimized baseline---not inflated by solver inefficiencies.
